\documentclass{article}
\usepackage{diagbox}
\usepackage{enumitem}
\usepackage{amssymb}

% Preamble
\title{Ejercicios de optimización 1}
\author{Bastián Ibañez Martíenz}
\date{24/03/2026}

\begin{document}
\maketitle
\newpage

\section*{2.7 Problema de transporte}

Daar S. A. es un fabricante de teléfonos móviles de MioXi de última
tecnología. La empresa tiene dos plantas de producción, una en Cali y otra
en Medellín, además, tiene un almacén ubicado en Bogotá. Daar vende
sus teléfonos en seis mercados: Pasto (PSO), Medellín (MDE), Cartagena
(CTG), Cali (CLO), Manizales (MZL) y Bogotá (Btá). \bigskip

Para satisfacer la demanda semanal del mercado, la gerente de distribución
debe determinar un plan para el envío desde las plantas y su almacén
hacia los distintos mercados, al mínimo costo. Ella tiene conocimiento de
que existen 600 unidades de MioXi disponibles para ser enviados desde la
planta de Cali, 600 desde la planta de Medellín y 250 desde el almacén
en Bogotá. El plan de envío no puede sobrepasar estos valores desde las
plantas y el almacén. \bigskip

Las demandas del mercado se indican en la Tabla 2.4, junto con el costo de
transporte asociado al envío de una unidad de teléfono MioXi. \bigskip

\begin{tabular}{|c|cccccc|}
\hline
& \multicolumn{6}{c|}{Destinos} \\
\cline{2-7}
Origen & PSO & MDE & CTG & CLO & MZL & BTA \\
\hline
Planta 1 Cali & 13150 & 9650 & 19950 & 5000 & 9650 & 15950 \\
Planta 2 Medellín & 19950 & 6500 & 19950 & 9650 & 11150 & 15950 \\
Almacén Bogotá & 17150 & 15950 & 19950 & 15950 & 15950 & 7000 \\
\hline
Demanda & 280 & 250 & 240 & 170 & 210 & 300 \\
\hline
\end{tabular} \bigskip

\begin{enumerate}[label=\alph*.]
    \item Verificar la viabilidad del problema de transporte de Daar, es decir,
¿la oferta será suficiente para satisfacer la demanda total?
    \item Formular un modelo matemático que asegure un costo de transporte
total mínimo al realizar los envíos de teléfonos MioXi desde las
plantas y el almacén, hacia los mercados.
\end{enumerate}

\subsection*{Definición de conjuntos}
\begin{itemize}
    \item I: Nodos de origen \{1,2,3\}
    \item J: Mercados \{1,2,3,4,5,6\}
\end{itemize}

\subsection*{Definición de Parámetros}
\begin{itemize}
    \item $c_{ij}$: Costo de transporter un teléfono desde el Origen $i \in I$ al mercado $j \in J$
    \item $o_{i}$: Oferta disponible en el origen $i \in I$
    \item $d_{i}$: Demanda en el mercado $j \in J$
\end{itemize}

\subsection*{Variables de decisión}
\begin{itemize}
    \item $X_{ij}$: Unidades de teléfonos $X$ a enviar el Origen $i \in I$ al mercado $j \in J$
\end{itemize}

\subsection*{Modelo algebráico}

Minimizar $\sum\limits_{i \in I}\sum\limits_{j \in J} c_{ij}*X_{ij}$ \bigskip

Sujeto a:
\begin{enumerate}
    \item $\sum\limits_{j \in J} X_{ij} \leq o{ij}, \hfill \forall i \in I$
    \item $\sum\limits_{i \in I} X_{ij} \geq d{ij}, \hfill \forall j \in J$
    \item $X_{ij} \in \mathbb{Z^+} \cup \{0\}, \hfill \forall i \in I, \forall j \in J$ 
\end{enumerate} 

\newpage
\section*{2.9 Problema del restaurante italiano}

Un restaurante italiano cuenta con cinco especialidades i en su menú:
tallarines, ravioles, lasaña, corbatitas y fideos. Cada una genera distintas
utilidades ui en dólares y tiene una demanda máxima diaria de especialidades
di. Ambos son datos estimados mediante un estudio de mercado, que se
resume en la Tabla \bigskip

\begin{tabular}{|c|ccccc|}
\hline
Especialidades & Tallarines & Ravioles & Lasaña & Corbatitas & Fideos \\
\hline
Utilidad [USD] & 50 & 65 & 40 & 70 & 55 \\
Demanda [und] & 15 & 12 & 32 & 18 & 10 \\
\hline
\end{tabular}
\bigskip

Cada especialidad i puede estar compuesta por los siguientes j ingredientes:
tomate, cebolla, huevo y pimiento. Así, cada especialidad i requiere cantidades
diferentes de ingredientes j para su elaboración (aij), como muestra la Tabla 2.7.
\bigskip

\begin{tabular}{|c|cccc|}
\hline
& Tomate & Cebolla & Huevo & Pimiento \\
\hline
Tallarines & 1 & 1 & 0 & 1 \\
Ravioles & 4 & 2 & 0 & 1 \\
Lasaña & 2 & 1 & 1 & 2 \\
Corbatitas & 0 & 0 & 2 & 0 \\
Fideos & 1 & 1 & 0 & 1 \\
\hline
\end{tabular}

\subsection*{Definición de conjuntos}
\begin{itemize}
    \item I: Especialidades \{1,2,3,4,5\}
    \item J: Ingredientes frescos \{1,2,3,4\}
\end{itemize}

\subsection*{Definición de parámetros}
\begin{itemize}
    \item $u_{i}$: Utilidad de la especialidad $i \in I$
    \item $d_{i}$: Demanda de la especialidad $i \in I$
    \item $a_{ij}$: Cantidad de ingredientes $j \in J$ requeridos para preparar la especialidad $i \in I$
    \item $m_{j}$: Cantidad disponible del ingrediente $j \in J$
\end{itemize}

\subsection*{Definición de variables de decisión}
\begin{itemize}
    \item $X_{i}$: Numero de especialidad diaria $i \in I$ a producir
\end{itemize}

\subsection*{Modelo algebráico}
Maximizar z = $u_{i} * X_{i}$ \bigskip
Sujeto a: \bigskip
\begin{enumerate}
    \item $\sum\limits_{i \in I} X_{i} * a_{ij} < m_{j}, \hfill \forall j \in J$
    \item $X_{i} \leq d_{i}, \hfill \forall i \in I$
    \item $X_{i} \in \mathbb{Z^+} \cup \{0\}, \hfill \forall i \in I$
\end{enumerate}

\section*{2.11 Problema de compra de zapatos}
Zapato a tu medida es una empresa de zapatos que cuenta con un
presupuesto de 50 000 USD para un proyecto de innovación. Para ello,
acude a 2 proveedores: Nikis y Pinki. Debido a que es una compra al por
mayor, ambos proveedores venden los zapatos en paquetes (cajas) de varias
unidades, que pueden variar dependiendo del tipo de zapato. La empresa
desea comprar la misma cantidad de cajas a ambos proveedores, pero la
compra debe ser de mínimo 30 cajas en total. Adicionalmente, la empresa
busca que el total de pares de zapatos generados por la compra de las cajas
sea al menos de 250 pares. \bigskip

En la Tabla 2.13 se observan los datos de costos, número de pares
de zapatos por caja y la utilidad que generaría para cada paquete y
fabricante, respectivamente. \bigskip

\begin{tabular}{|c|cccc|}
\hline
Cajaj & Fabricante & Costo [USD] & Número de pares de zapatos & Utilidad por caja [USD/caja] \\
\hline
Tenis & Nikis & 2700 & 45 & 4500 \\
Botín & Nikis & 2400 & 30 & 2550 \\
Tacones & Pinki & 640 & 16 & 960 \\
Sandalias & Pinki & 625 & 25 & 1250 \\
\hline
\end{tabular}

¿Cuántas cajas de zapatos debe comprar para obtener la máxima utilidad?

\subsection*{Definición de conjuntos}
\begin{itemize}
    \item I: Tipo de caja \{1,2,3,4\}
\end{itemize}

\subsection*{Definición de parámetros}
\begin{itemize}
    \item $c_{i}$: Costo de la caja $i \in I$
    \item $o_{i}$: Número de pares de zapatos en la caja $i \in I$
    \item $u_{i}$: Utilidad por caja de la caja $i \in I$
\end{itemize}

\subsection*{Definición de variables de decisión}
\begin{itemize}
    \item $x_{i}$: Unidades a comprar de la caja $i \in I$
\end{itemize}

\subsection*{Modelo algebráico}

Maximizar: $\sum\limits_{i \in I} u_{i}*x_{i}$ \\
Sujeto a:
\begin{enumerate}
    \item $\sum\limits_{i \in I} x_i * c_i \leq 50000$
    \item $\sum\limits_{i \in I} x_i \geq 30$
    \item $\sum\limits_{i=1}^2 x_i = \sum\limits_{i=3}^4 x_i$
    \item $\sum\limits_{i \in I} x_i * o_i \geq 250$
    \item $x_i \in \mathbb{Z^+} \cup \{0\}, \hfill \forall i \in I$

\end{enumerate}

\newpage
\section*{3.3 Problema del inventario de tapabocas}

Una empresa distribuidora de tapabocas tiene la siguiente demanda
pronosticada para los próximos cuatro meses: 2 800, 2 200, 3 200 y 2 500 unidades
para los meses 1, 2, 3 y 4, respectivamente. La empresa tiene capacidad
para producir 2 700 unidades de tapabocas por mes, en sus turnos normales.
Utilizando tiempo extra, es posible fabricar 300 unidades adicionales. La
producción en tiempo extra tiene un sobrecosto de 10 USD por unidad. \bigskip
La administración ha estimado un costo de almacenamiento de 2 USD por
unidad producida y que no sea vendida en un mes determinado, como se
muestra a continuación: \\\\
\begin{center}
    $\frac{USD}{uni*mes} $
\end{center}
¿Cuál es el programa óptimo de producción que minimice los costos de
producción en tiempo extra y almacenamiento, cumpliendo la demanda de
tapabocas cada mes? Asumir que la cantidad en existencia actual es 0 y se
desea un inventario final del periodo también igual a 0.

\subsection*{Datos}
\begin{itemize}
    \item 4 meses
    \item 2700 capacidad normal
    \item 300 capacidad extra a \$10/unidad
    \item Costo de almacenamiento \$2/unidad no vendida
\end{itemize}

\subsection*{Definición de conjuntos}
\begin{itemize}
    \item T: Mes de producción \{1,2,3,4\}
\end{itemize}

\subsection*{Definición de parámetros}
\begin{itemize}
    \item $d_t$: Demanda proyectada para el mes $t \in T$
    \item $h$: Costo de inventario
    \item $p_e$: Costo de producción de tapabocas en tiempo extra
    \item $cap$: Capacidad de producción en tiempo regular
    \item $cap_e$: Capacidad de producción en tiempo extra
\end{itemize}

\subsection*{Definición de variables de decisión}
\begin{itemize}
    \item $x_t$: Unidades a producir en tiempo regular, $t\in T$
    \item $e_t$: Unidades a producir en tiempo extra, $t\in T$
    \item $i_t$: Unidades a tener en inventario, $t\in T \cup \{0\}$ 
\end{itemize}

\subsection*{Restricciones}
\begin{itemize}
    \item $x_t \leq cap$
    \item $e_t \leq cap_e$
    \item $x_t + e_t \geq d_t$  
\end{itemize}

\subsection*{Modelo algebráico}
Minimizar z: $\sum\limits_{t \in T} i_t*h + p_e * e_t$ \\
Sujeto a: \\
$x_t \leq cap, \hfill \forall t \in T$ \\
$e_t \leq cap_e, \hfill \forall t \in T$ \\
$i_{t-1} + x_t + e_t \geq d_t, \hfill \forall t \in T$ \\
$x_t, e_t, \in \mathbb{Z^+} \cup \{0\}$ \\


\end{document}
